\documentclass[twoside]{article}
\usepackage[table]{xcolor}
\usepackage[paper=letterpaper,left=3cm,right=3cm,top=3cm,bottom=2cm,includefoot]{geometry}
\usepackage{fancyhdr}
\usepackage{tabularx}
\usepackage{multirow}
\usepackage[colorlinks]{hyperref}
\usepackage{hhline}
\usepackage{amsmath}
\usepackage{enumitem}
\setenumerate{itemsep=-3pt,topsep=3pt}
\setdescription{itemsep=-3pt,topsep=3pt,leftmargin=!,labelwidth=4.5cm}
\usepackage{pgfgantt}
\usepackage{graphicx}   \graphicspath{{./img/} {./tex/img/}}
\usepackage{amsfonts}
\usepackage[spanish]{babel}
\usepackage[utf8]{inputenc}
\usepackage[T1]{fontenc}
\usepackage{listings}
\usepackage{xcolor}
\usepackage[most]{tcolorbox}
\usepackage{lipsum}
\def\code#1{\texttt{#1}}


% insertar codigo Julia en documento latex
%\usepackage{minted}
%\definecolor{LightGray}{gray}{0.92}

% Se define el entorno para Julia con las mismas configuraciones
%\newminted{julia}{ linenos,breaklines,mathescape,texcomments,xleftmargin=\parindent, numbersep=8pt, bgcolor=LightGray
%}

%\renewcommand\listingscaption{C\'odigo}


% bibliografia: descomente estas dos lineas para usar estilo numerico, e.g. [1].
\usepackage[square,numbers,sort&compress]{natbib}
\bibliographystyle{apastyle}

\pagestyle{fancy}
\renewcommand{\footrulewidth}{0.4pt}
\renewcommand{\headrulewidth}{0.4pt}
\fancyfoot{}

%header y footer
\fancyfoot[RE,RO]{\thepage}
\fancyfoot[LE,LO]{Benjamin Ayala Baeza}


\definecolor{tcc}{RGB}{217,217,217} % Table cell color

\renewcommand\tabularxcolumn[1]{m{#1}}
\setlength{\arrayrulewidth}{0.5pt}
\renewcommand{\arraystretch}{2}

\tcbset{
    portada/.style={
        enhanced,
        colback=white,
        colframe=black,
        sharp corners,
        boxrule=0.4pt,
        drop shadow, % sombra elegante
        width=\textwidth,
        center
    }
}

%esto hace que desaperezca el header en la primera plana
\fancypagestyle{plainfirst}{
    \fancyhf{}
    \fancyfoot[RE,RO]{\thepage} % mantiene número de página
}

\begin{document}

\thispagestyle{plainfirst}

%\vspace*{\fill}
\begin{tcolorbox}[portada]
    \centering
    {\Huge \bfseries Electric Field of a Point Charge}\\[0.5cm]
    %{\LARGE Derivación Númerica}\\[0.5cm]
    {\large Benjamín Ayala Baeza}
\end{tcolorbox}
%\vspace*{\fill}



\section{Introduction}
An usual problem in physics, specifically in electromagnetism is the finding of the electric field in a given situation,
Normally, we go from the most basic situation into the mot difficult ones. That is why I'm begging with most simple and intuitive 
search for the electric field, the field of just one charged particle in the space. I'll also compute it's electric potential. 

This is an interesting problem because it cements the foundations for more non-trivial situations. 
A simulation will be included, which shows how the electric field changes according a change in the charge's position.

\section{Physical Theory}
The most basic and direct way to compute the electric field of a distribution charge (continuos) is using Gauss's Law:

    \begin{align}
        \oint \vec{E} \cdot d\vec{S} = \frac{\rho}{\varepsilon_0}.  \label{eq:gauss-law}
    \end{align}
In general, solving this integral requires integrating over a known space but in general, we write the electric field of a distribution charge as:

    \begin{align}
        \vec{E}(\vec{r}) = \frac{1}{4\pi \varepsilon_0} \frac{\rho}{|\vec{r}-\vec{r}'|} \hat{r}. \label{eq:electric-field}
    \end{align}
Notices that if we are integrating over an enclosed sphere in \eqref{eq:gauss-law} we obtain the equation \eqref{eq:electric-field} in which we defined two vector, $\vec{r}$ and $\vec{r}'$,
where the first corresponds to the position of the field point whereas the second vector corresponds to the position of the charge distribution.
\end{document}
